\documentclass{ximera}
\title{Logarithmic Functions CW2}

\newcommand{\pskip}{\vskip 0.1 in}

\begin{document}
\begin{abstract}
Applications of Logarithmic functions.
\end{abstract}
\maketitle



\begin{question} \label{QPfrer3rdefs}
The function
\[
  C = g(t) = 20 + 80 \cdot  2^{-\frac{t}{4}} \, , \, t \geq 0,
\]
expresses the temperature (in Celsius degrees) of a cup of coffee in terms of the number of minutes since noon. The coffee is in a room held at a constant temperature of $20^\circ$ C.

\begin{enumerate}
\item Find the temperature of the coffee at 12:12pm without a calculator.

\item Describe the action of $g$ on its input.

\item Describe the action of $g^{-1}$ on its input.

\item Use your description in part (c) to find when the temperature of the coffee is $60^\circ$C. Give an exact time \emph{without} a calculator. Then use a calculator to approximate the time to the nearest minute.

\item Use your description in part (c) to find an expression for the function $g^{-1}$. Use the appropriate variable names for the input and output. Include a domain.

\item When is the temperature of the coffee $20^\circ$C?

\item Use algebra to show $g^{-1}$ undoes $g$.

\item How long does it take the temperature difference between the coffee and the room to halve?
\end{enumerate}

\end{question}


\begin{question}  \label{Q:dfehy3190h5h}
The balance in an account grows exponentially. The account takes $5$ years to grow from $\$10,000$ to $\$15,000$. 

\begin{enumerate}
\item Describe how the balance grows. Be precise.

\item Find a function 
\[
B = f(t), -5\leq t \leq 20,
\]
that expresses the balance in the account (in dollars) in terms of the number of years since the balance was $\$10,000$.

\item Enter your function in Line 1 of the worksheet below. 

\begin{onlineOnly}
    \begin{center}
\desmos{depmu5ngz0}{900}{600}
\end{center}
\end{onlineOnly}


\href{https://www.desmos.com/calculator/depmu5ngz0}{141: Compound Interest 1}

\item Use your graph from part (c) to approximate how long it takes the balance to increase from $\$15,000$ to $\$20,000$. Then compute the exact time \emph{without} a calculator. Then give a decimal approximation to the nearest hundredth of a year.

\item Use your graph from part (c) to approximate how long it takes the balance to increase from $\$40,000$ to $\$45,000$. Then compute the exact time \emph{without} a calculator. Then give a decimal approximation to the nearest hundredth of a year.

\item You should have found that
\[
    B = f(t) = 10000 (1.5)^{\frac{1}{5}t} , -5 \leq t \leq 20.
\]

\begin{enumerate}

\item Describe how the function $f$ acts on its input.

\item Use your description above to describe the action of $f^{-1}$.

\item Use your description of how $f^{-1}$ acts on its input to find an expression for the function 
\[ 
      t  = f^{-1}(B) .
\]

\item Use algebra and solve the equation $B=f(t)$ for $t$ to find an expression for the function $t = f^{-1})(B)$. Check that you get the same result.

\item State the domain and range of $f^{-1}$ using set-buider notation. Describe also the meaning of this function. What does it take as an input? What does it return as an output.

\item Use the graph of $B=f(t)$ above to sketch \emph{by hand} a graph of $t=f^{-1}(B)$.

\end{enumerate}

\item Use the graph above to estimate the time it takes the balance to increase by $25\%$.

\item Use your expression for $f^{-1}$ to determine the exact time it takes the balance to increase by $25\%$. Then use a calculator to approximate this time and compare this with your estimate.

\end{enumerate}
\end{question}


\end{document}